% Consolidated 25-elastic-waves from 0d2a41587ad18cded191ec1d80e8718df37ec2fe; original TeX fragments preserved.

% === BEGIN 25-elastic-waves-150-165.tex blob 1c0581980ff97154ce279eb05d674224fdbaefa2 ===
% Irodov Part IV §4.3 original 150--165.
% Source user PDF pp.184--187 (printed pp.175--178), answers pp.320--322.
\waveproblem{150}{175--176}
空气中 A、B 两点相距 $\ell$，
其间绝对温度沿直线由 $T_1$ 线性变化到 $T_2$。
声速与温度的关系为
$v(T)=\alpha\sqrt T$，$\alpha$ 为常数。
求声波从 A 到 B 所需时间。
\exsolution
令 $x=0$ 为 A，$x=\ell$ 为 B，
$T(x)=T_1+(T_2-T_1)x/\ell$。
传播时间是局部速度倒数的路径积分：
\[
t=\int_0^\ell\frac{\mathrm dx}{\alpha\sqrt{T(x)}}
=\frac{2\ell}{\alpha}
\frac{\sqrt{T_2}-\sqrt{T_1}}{T_2-T_1}.
\]
因
$T_2-T_1=(\sqrt{T_2}-\sqrt{T_1})
(\sqrt{T_2}+\sqrt{T_1})$，得
\[
\boxed{t=\frac{2\ell}
{\alpha(\sqrt{T_1}+\sqrt{T_2})}}.
\]
\excheck 当 $T_2\to T_1=T$ 时，
$t\to\ell/(\alpha\sqrt T)$；
底本 PDF 文字抽取把 $v=\alpha\sqrt T$
中的根号部分遗漏，已据物理单位与书末恢复。

\waveproblem{151}{176}
角频率为 $\omega$、相速度为 $v$ 的平面谐波，
传播方向与 $x,y,z$ 正轴分别成角
$\alpha,\beta,\gamma$。
求介质两点
$\bm r_1=(x_1,y_1,z_1)$、
$\bm r_2=(x_2,y_2,z_2)$ 之间的振动相位差。
\exsolution
传播方向单位矢量
$\bm n=(\cos\alpha,\cos\beta,\cos\gamma)$，
波矢 $\bm k=(\omega/v)\bm n$。
波动相位
$\phi(\bm r,t)=\omega t-\bm k\cdot\bm r+\phi_0$。
同一时刻两点相位差
$\Delta\phi=-\bm k\cdot(\bm r_2-\bm r_1)$。
因此
\[
\boxed{|\Delta\phi|=\frac{\omega}{v}
\left|
(x_2-x_1)\cos\alpha+
(y_2-y_1)\cos\beta+
(z_2-z_1)\cos\gamma
\right|}.
\]
\excheck 必须满足方向余弦恒等式
$\cos^2\alpha+\cos^2\beta+\cos^2\gamma=1$。

\waveproblem{152}{176}
一列角频率为 $\omega$ 的平面波，
同一相位前沿沿三坐标轴分别以速度
$v_1,v_2,v_3$ 经过固定坐标。
用 $\bm e_x,\bm e_y,\bm e_z$ 表示波矢。
\exsolution
相位面满足
$\omega t-k_xx-k_yy-k_zz=\text{const}$。
沿 x 轴取 y,z 固定，恒相位点的速度
$\mathrm dx/\mathrm dt=\omega/k_x=v_1$，
故 $k_x=\omega/v_1$；
其它方向同理。
于是
\[
\boxed{\bm k=\omega\left(
\frac{\bm e_x}{v_1}+
\frac{\bm e_y}{v_2}+
\frac{\bm e_z}{v_3}\right)}.
\]
\excheck 各 $v_i$ 是波面在\emph{坐标轴方向的相位交点速度}，
并非波速矢量的三个分量；
可有 $v_i>|\bm v|$ 而不违背波速定义。

\waveproblem{153}{176}
介质静止参考系 K 内，向 +x 传播的平面弹性波
$\xi=a\cos(\omega t-kx)$。
参考系 K' 以速度 $V$ 沿 +x 相对介质运动。
求 K' 中的波方程，讨论观测频率与传播方向。
\exsolution
非相对论介质中的坐标变换
$x=x'+Vt$、$t'=t$。
将其代入原相位：
\[
\boxed{\xi'(x',t')=
a\cos[(\omega-kV)t'-kx']}.
\]
观测有向角频率为
$\omega'=\omega-kV=\omega(1-V/v)$，
其中原波相速度 $v=\omega/k$；
相对 K' 的有向相速度
\[
\boxed{v'=v-V}.
\]
当 $V=v$，观察者随波面同行，场的时域变化静止；
$V>v$ 时有向相位变化反号，
观测到波面相对 K' 沿 -x 运动。
\excheck 这是弹性波在介质中的伽利略坐标变换，
不能无条件替换为真空光波的相对论变换。

\waveproblem{154}{176}
证明任何二阶可微函数
$\xi=f(t+\alpha x)$
在适当 $\alpha$ 下满足一维波动方程，并说明 $\alpha$ 的物理意义。
\exsolution
令 $u=t+\alpha x$，
链式法则给
\[
\frac{\partial^2\xi}{\partial t^2}=f''(u),
\qquad
\frac{\partial^2\xi}{\partial x^2}=\alpha^2f''(u).
\]
一维波动方程
$\partial_x^2\xi=v^{-2}\partial_t^2\xi$
要求
$(\alpha^2-v^{-2})f''(u)=0$。
对任意非平凡 $f$，必须有
\[
\boxed{\alpha=\pm\frac1v}.
\]
若 $\alpha=-1/v$，
相位 $t-x/v=\text{const}$ 对应沿 +x 传播；
若 $\alpha=+1/v$ 则沿 -x。
\excheck 题干仅说“differentiable”略弱于含二阶导数的经典波动方程要求；
作经典解至少需要二阶可微，弱解则另论。

\waveproblem{155}{176}
行波
$\xi=60\cos(1800t-5.3x)$，
其中 $\xi$ 单位为微米、$t$ 用秒、$x$ 用米。
求：(a) 介质质点位移振幅与波长之比；
(b) 质点振动速度幅度及其与传播速度之比；
(c) 纵向相对形变的振幅及其与质点速度幅度的关系。
\exsolution
振幅 $a=60\times10^{-6}$m，
角频率 $\omega=1800\,\mathrm{s^{-1}}$，
波数 $k=5.3\,\mathrm{m^{-1}}$。
波长 $\lambda=2\pi/k$，
传播速度 $v=\omega/k\simeq340$m/s。
(a)
\[
\boxed{\frac a\lambda=\frac{ak}{2\pi}
\simeq5.1\times10^{-5}}.
\]
(b) 质点速度幅度
\[
\boxed{v_m=a\omega=0.108\,\mathrm{m/s}
\simeq11\,\mathrm{cm/s}},
\quad
\boxed{\frac{v_m}{v}=\frac{a\omega}{\omega/k}
=ak\simeq3.2\times10^{-4}}.
\]
(c) 相对形变
$\epsilon=\partial_x\xi
=ak\sin(\omega t-kx)$，
因此
\[
\boxed{\epsilon_m=ak\simeq3.2\times10^{-4}},
\quad
\boxed{v_m=v\epsilon_m}.
\]
\excheck 传播速度和质点速度不同；
质点速度约0.1m/s，波传播约340m/s。

\waveproblem{156}{176}
平面波
$\xi=a\cos(\omega t-kx)$
在均匀弹性介质传播。
于 $t=0$ 绘出
(a) $\xi$、$\partial_t\xi$、$\partial_x\xi$ 随 $x$ 的曲线；
(b) $\xi=0$ 处质点速度方向（纵波和横波）；
(c) 纵波中密度变化曲线。
\exsolution
直接求导并令 $t=0$：
\[
\boxed{\xi(x,0)=a\cos kx},\quad
\boxed{\partial_t\xi(x,0)=a\omega\sin kx},\quad
\boxed{\partial_x\xi(x,0)=-ak\sin kx}.
\]
因此两导数曲线均是以 $kx$ 为相位的正弦波，
彼此反相，而位移曲线是余弦波。
在零位移点 $kx=(n+1/2)\pi$，
速度模达到最大；
纵波中质点速度沿传播轴，其正负由
$\sin kx$ 决定；
横波中质点速度沿指定的横向偏振轴，同样由
$\sin kx$ 控制。
连续性方程的一阶近似
$\delta\rho/\rho_0=-\partial_x\xi$，
故
\[
\boxed{\rho(x,0)\simeq
\rho_0[1+ak\sin kx]}.
\]
\excheck 密度极大处对应局部压缩
$\partial_x\xi<0$，与质量守恒一致；
该三组函数即可逐点绘制原题所求曲线。

\waveproblem{157}{176}
均匀介质中衰减平面波
$\xi=ae^{-\gamma x}\cos(\omega t-kx)$。
两点处质点位移振幅相差
$\eta=1.0\%$（按前一点振幅减少百分比），
给定 $\gamma=0.42\,\mathrm{m^{-1}}$、
波长 $\lambda=50$cm。
求两点的相位差。
\exsolution
若位置差为 $\Delta x>0$，
后一点振幅
$A_2=A_1e^{-\gamma\Delta x}
=(1-\eta)A_1$，
则
\[
\Delta x=-\frac{\ln(1-\eta)}{\gamma}.
\]
相位差模
$|\Delta\phi|=k\Delta x$，
$k=2\pi/\lambda$，
故
\[
\boxed{|\Delta\phi|=
-\frac{2\pi}{\gamma\lambda}\ln(1-\eta)}.
\]
代 $\eta=0.010$：
\[
\boxed{|\Delta\phi|\simeq0.30\,\mathrm{rad}}.
\]
\excheck 使用 $\ln(1-\eta)\approx-\eta$ 时近似结果
$2\pi\eta/(\gamma\lambda)$ 同样约0.30。

\waveproblem{158}{176--177}
均匀无吸收介质中的点球面波源 P
已知位于两位置矢量分别为
$\bm r_1,\bm r_2$ 的点之间的直线上。
该两点处质点振动幅度分别为 $a_1,a_2$。
求源点位置矢量。
\exsolution
无吸收的球面波位移振幅
$a(r)=A/r$，
其中 $r$ 是离声源的距离。
设两点到源的距离为 $d_1,d_2$，
则
$a_1d_1=a_2d_2=A$。
因为源在两点之间，
$d_1+d_2=d=|\bm r_2-\bm r_1|$。
解得
$d_1=a_2d/(a_1+a_2)$、
$d_2=a_1d/(a_1+a_2)$。
源点相当于从 $\bm r_1$ 沿两点连线走比例
$d_1/d$，
所以
\[
\boxed{\bm r_P=
\frac{a_1\bm r_1+a_2\bm r_2}{a_1+a_2}}.
\]
\excheck 当 $a_1>a_2$，
源更靠近点1，因此位置权重应使结果更靠近 $\bm r_1$。

\waveproblem{159}{177}
频率 $\nu=1.45$kHz 的各向同性点声源，
距离 $r_0=5.0$m 处质点位移振幅
$a_0=50\,\mathrm{\mu m}$。
距声源 $r=10.0$m 的 A 点振幅比 $a_0$ 小
$\eta=3.0$ 倍。
求：(a) 振幅衰减系数 $\gamma$；
(b) A 点质点速度幅度。
\exsolution
衰减球面波幅度随径向距离变化
$a(r)=a_0(r_0/r)e^{-\gamma(r-r_0)}$。
由 $a_0/a(r)=\eta$，
\[
\eta=\frac r{r_0}e^{\gamma(r-r_0)},
\qquad
\boxed{\gamma=\frac{\ln(\eta r_0/r)}{r-r_0}}.
\]
代 $\eta r_0/r=1.5$、$r-r_0=5$m，
\[
\boxed{\gamma\simeq0.081\,\mathrm{m^{-1}}}.
\]
A 点位移幅度 $a(r)=a_0/\eta$，
质点速度幅度
\[
\boxed{v_m=2\pi\nu\,\frac{a_0}{\eta}
\simeq0.15\,\mathrm{m/s}=15\,\mathrm{cm/s}}.
\]
\excheck 既要包括球面波几何扩散 $r_0/r$，
也要包括材料吸收 $e^{-\gamma(r-r_0)}$；不能只用指数衰减。

\waveproblem{160}{177}
同种弹性介质中，两列同频同幅平面波分别沿 x、y 轴传播，
振动分别为
$\xi_1=a\cos(\omega t-kx)$、
$\xi_2=a\cos(\omega t-ky)$。
求 $xy$ 平面各点质点的合运动情况：
(a) 两波均为振动方向相同的横波；
(b) 两波均为各自沿传播方向振动的纵波。
\exsolution
(a) 合位移沿共同横向方向：
\[
\xi=\xi_1+\xi_2
=2a\cos\frac{k(y-x)}2
\cos\left(\omega t-\frac{k(x+y)}2\right).
\]
故
$\boxed{y-x=n\lambda}$ 为振幅极大线，
$\boxed{y-x=(n+\tfrac12)\lambda}$ 为结点线，
$n\in\mathbb Z$。
(b) 对纵波，
质点的二维位移分量为
$x_{\rm osc}=a\cos(\omega t-kx)$、
$y_{\rm osc}=a\cos(\omega t-ky)$。
两分量同频、相位差
$\delta=k(y-x)$：
当 $\delta=n\pi$ 时为直线运动，
其中偶数 $n$ 沿 $(1,1)$、奇数 $n$ 沿 $(1,-1)$；
当 $\delta=\pi/2+n\pi$ 时为圆周运动，
其它情形一般为椭圆。
\excheck 两波的\emph{纵波}位移沿不同方向，
不能标量相加代替二维参数轨迹。

\waveproblem{161}{177}
均匀介质中无衰减的平面谐波传播。
某固定点在经过质点位移最大时刻后的
$T/6$，测得\emph{瞬时总振动能量密度}为 $w_0$。
求该波的时间平均总能量密度。
\exsolution
平面行波的动能与弹性势能密度相同，
其总能量密度随相位
$\theta$ 为
$w(\theta)=\rho a^2\omega^2\sin^2\theta$。
位移最大时 $\theta=0$，
再过 $T/6$：
$\theta=2\pi(T/6)/T=\pi/3$，
所以
$w_0=\rho a^2\omega^2\sin^2(\pi/3)
=\frac34\rho a^2\omega^2$。
时间平均
$\overline w=\frac12\rho a^2\omega^2$，
从而
\[
\boxed{\overline w=\frac23w_0}.
\]
\excheck 这是某时刻的\emph{总}能量密度与周期平均的关系，
不应把动能或势能单项误当总能。

\waveproblem{162}{177}
各向同性点声源位于圆环所在平面法线上的一点，
离圆心 O 距离 $\ell=1.00$m，
圆环半径 $R=0.50$m。
已知 O 点平均声强
$I_0=30\,\mathrm{\mu W/m^2}$，
忽略衰减，求通过圆环所围圆盘面的平均能流。
\exsolution
由球面扩散，
声源总功率
$P=4\pi\ell^2I_0$。
圆盘从声源所张立体角
$\Omega=2\pi[1-\cos\theta]$，
$\cos\theta=\ell/\sqrt{\ell^2+R^2}$。
穿过圆盘的功率为
$P\Omega/(4\pi)$，故
\[
\boxed{\Phi=
2\pi\ell^2I_0
\left[1-\frac{\ell}{\sqrt{\ell^2+R^2}}\right]}.
\]
代数值
$\boxed{\Phi\simeq2.0\times10^{-5}\,\mathrm W
=20\,\mathrm{\mu W}}$。
\excheck 必须计算声流的\emph{法向投影}，不能只以圆盘面积乘中心声强，
因为声强沿圆盘位置及方向均变化。

\waveproblem{163}{177}
各向同性点声源功率 $P=0.10$W，
位于空心圆柱中心，
圆柱半径 $R=1.0$m、高度 $h=2.0$m。
假定圆柱壁完全吸收声波，求达到圆柱\emph{侧面}的平均能流。
\exsolution
对从源点出发的射线，
若与轴的极角为 $\theta$，
当 $|\tan\theta|>2R/h$ 时先撞侧壁，
否则先撞端盖。
因此两端盖张的两个锥形区域所占功率比例为
$1-\cos\theta_0$，
其中
$\cos\theta_0=(h/2)/\sqrt{R^2+(h/2)^2}$。
侧壁所得比例为 $\cos\theta_0$，故
\[
\boxed{\Phi_{\rm side}
=P\frac{h/2}{\sqrt{R^2+(h/2)^2}}
=\frac P{\sqrt{1+(2R/h)^2}}}.
\]
代 $h=2R$，
\[
\boxed{\Phi_{\rm side}\simeq0.071\,\mathrm W}.
\]
\excheck 只有部分球面立体角对应侧壁，
不能把总功率 $P$ 全部归入侧面。

\waveproblem{164}{177}
均匀介质内平面驻波
$\xi=a\cos kx\,\cos\omega t$。
作出：(a) $t=0$ 与 $t=T/2$ 的位移和空间导数；
(b) 同时刻的纵波介质密度分布；
(c) $t=T/4$ 质点速度分布及各腹点速度方向。
\exsolution
位移、空间导数与质点速度为
\[
\xi=a\cos kx\cos\omega t,\quad
\frac{\partial\xi}{\partial x}
=-ak\sin kx\cos\omega t,\quad
u=\frac{\partial\xi}{\partial t}
=-a\omega\cos kx\sin\omega t.
\]
在 $t=0$：
$\xi=a\cos kx$、
$\partial_x\xi=-ak\sin kx$；
在 $t=T/2$：
$\xi=-a\cos kx$、
$\partial_x\xi=+ak\sin kx$。
纵波一阶密度
$\rho(x,t)\simeq\rho_0[1-\partial_x\xi]$，故
\[
\boxed{\rho(x,0)=\rho_0(1+ak\sin kx)},\quad
\boxed{\rho(x,T/2)=\rho_0(1-ak\sin kx)}.
\]
在 $t=T/4$，
$\boxed{u(x,T/4)=-a\omega\cos kx}$。
腹点 $kx=n\pi$ 的质点速度交替反向，
纵波时沿 x、横波时沿横向偏振轴。
\excheck 位移结点在 $\cos kx=0$，速度始终为零；
邻近腹点的质点速度在该时刻方向相反。

\waveproblem{165}{177--178}
均匀介质密度为 $\rho$，
纵向平面驻波
$\xi=a\cos kx\cos\omega t$。
求：(a) 弹性势能密度；
(b) 动能密度；
并给出在相邻两个位移结点之间，
$t=0$、$t=T/4$ 的\emph{总能量密度}分布。
\exsolution
介质纵向线性应变
$\epsilon=\partial_x\xi=-ak\sin kx\cos\omega t$。
弹性势能密度为
$\tfrac12\rho v^2\epsilon^2$，
利用 $\omega=vk$：
\[
\boxed{w_p(x,t)=\frac12\rho a^2\omega^2
\sin^2kx\,\cos^2\omega t}.
\]
质点速度
$\dot\xi=-a\omega\cos kx\sin\omega t$，
所以
\[
\boxed{w_k(x,t)=\frac12\rho a^2\omega^2
\cos^2kx\,\sin^2\omega t}.
\]
$t=0$ 时动能为零，
$\boxed{w(x,0)=\tfrac12\rho a^2\omega^2\sin^2kx}$；
$t=T/4$ 时势能为零，
$\boxed{w(x,T/4)=\tfrac12\rho a^2\omega^2\cos^2kx}$。
\excheck 两曲线互为四分之一波长的空间相位平移；
相邻结点间的\emph{积分总能量}不变，尽管局部能量密度随时间迁移。
% === END 25-elastic-waves-150-165.tex ===

% === BEGIN 25-elastic-waves-166-181.tex blob b1a2d4875ce2eff417774c69770aa6c39c071d99 ===
% Irodov Part IV §4.3 original 166--181. User PDF pp.187--188
% printed pp.178--179, answer PDF pp.322--323 (printed pp.321--323).
\waveproblem{166}{178}
长 $L=120$cm 的弦两端固定并形成驻波。
在弦上位移\emph{振幅}等于
$3.5$mm 的相邻点之间相距
$d=15.0$cm。
求弦振幅最大值，并指出所属泛音。
\exsolution
驻波振幅空间分布
$A(x)=A_{\max}|\sin kx|$。
方程 $|\sin kx|=3.5/A_{\max}$
在每个半波长周期内通常有两组解，相邻解点间隔交替。
题设把它们描述为相邻点间距均为同一个 $d$，
因此两类间距相等，解的相位必须处于
$|\sin kx|=1/\sqrt2$。
于是
\[
\boxed{A_{\max}=3.5\sqrt2\,\mathrm{mm}
\simeq5.0\,\mathrm{mm}}.
\]
这种情况下相邻指定点的相位差为 $\pi/2$，
$d=\lambda/4$，故 $\lambda=4d=0.60$m。
两端固定有
$L=n\lambda/2$，所以
$n=2L/\lambda=4$，
即第四谐波，
\[
\boxed{\text{第三泛音（third overtone）}}.
\]
\excheck “第三泛音”对应第四谐波而非第三谐波；
且数值 $5.0$mm 与原版答案一致。

\waveproblem{167}{178}
两根最初完全相同的弦，
分别被拉长
$\eta_1=2.0\%$、
$\eta_2=4.0\%$。
认为张力与相对伸长成正比。
求两根弦的基频之比。
\exsolution
设原弦长 $L_0$、原线密度 $\mu_0$、
材料张力常数为 $K$。
伸长比例 $\eta$ 后长度
$L=L_0(1+\eta)$，
张力 $T=K\eta$，
整根弦的质量不变，
线密度 $\mu=\mu_0/(1+\eta)$。
两端固定弦基频
$f=(1/(2L))\sqrt{T/\mu}$，
所以
$f\propto\sqrt{\eta/(1+\eta)}$。
因此
\[
\boxed{\frac{f_2}{f_1}
=\sqrt{\frac{\eta_2(1+\eta_1)}
{\eta_1(1+\eta_2)}}}
=\sqrt{\frac{0.04(1.02)}{0.02(1.04)}}
\simeq1.4.
\]
\excheck 此题不只是 $f\propto\sqrt T$，
因为伸长还改变振动部分的弦长及线密度。

\waveproblem{168}{178}
一根均匀张紧金属丝被缩短 $35\%$，
同时张力增加 $70\%$。
若弦材料的线密度保持不变，
求基频改变的方向及倍数。
\exsolution
弦两端固定基频
$f=(1/(2L))\sqrt{T/\mu}$，
缩短后
$L'=0.65L$，
新张力
$T'=1.70T$。
若取同材均匀金属丝重新截短而线密度不变，
\[
\boxed{\frac{f'}f=\frac1{0.65}\sqrt{1.70}
\simeq2.0}.
\]
因此基频\emph{升高至约2倍}。
\excheck 题目中“shortened”按截短同种金属丝理解；
如果是将原丝整体机械压短而改变密度/线密度，则还需额外模型。

\waveproblem{169}{178}
用一端为振动膜、另一端为可滑动活塞的声学共振管，
发声频率 $\nu=2000$Hz。
测得活塞的\emph{相邻共振位置}间距
$d=8.5$cm，求空气中的声速。
\exsolution
在相同端部边界条件下，
管长每增加半个波长，
便出现相邻一个纵向驻波共振：
$d=\lambda/2$。
故
$\lambda=2d=0.17$m，
声速
\[
\boxed{v=\nu\lambda=2\nu d
=2(2000)(0.085)
=340\,\mathrm{m/s}}.
\]
\excheck 不必知道端部修正的具体数值，
因为相邻共振位置的差将其消去。

\waveproblem{170}{178}
长 $L=85$cm 的空气柱中声速
$v=340$m/s。
统计低于
$\nu_{\max}=1250$Hz 的本征振动数：
(a) 一端封闭、一端开放；
(b) 两端均开放。
开放端视为位移腹点。
\exsolution
(a) 一端固定为位移结点，
另一端为位移腹点，
仅允许奇数四分之一波长：
\[
\boxed{\nu_n^{(a)}
=\frac{(2n+1)v}{4L}
=(2n+1)\,100\,\mathrm{Hz}},
\quad n=0,1,\ldots.
\]
小于1250Hz的为
100、300、500、700、900、1100Hz，
共
$\boxed{6\text{个}}$。
(b) 两端都是位移腹点：
\[
\boxed{\nu_n^{(b)}
=\frac{(n+1)v}{2L}
=(n+1)\,200\,\mathrm{Hz}},
\quad n=0,1,\ldots.
\]
小于1250Hz的为
200、400、600、800、1000、1200Hz，
也是
$\boxed{6\text{个}}$。
\excheck 频率上限是\emph{低于}1250Hz，
不包含等于上限的模态。

\waveproblem{171}{178}
铜棒长度 $L=50$cm，中点被夹紧，
求频率在20—50kHz之间的\emph{纵向}固有振动模式数及各个频率。
\exsolution
一维纵波声速
$v=\sqrt{E/\rho}$。
取铜的杨氏模量约
$E\simeq1.28\times10^{11}\,\mathrm{Pa}$、
密度约
$\rho=8.9\times10^3\,\mathrm{kg/m^3}$，
约得 $v\simeq3.79\,\mathrm{km/s}$，下文按书末约 $3.8\,\mathrm{km/s}$ 计算。铜的有效纵波速度依物性与试样而变，原题未给精确常数。
杆中点是位移结点，
两自由端为应力零、位移腹点；
所以每半根杆只能容纳奇数四分之一波长，
本征频率
\[
\boxed{\nu_n=\frac{(2n+1)v}{2L}
\simeq3.8(2n+1)\,\mathrm{kHz}},
\quad n=0,1,\ldots.
\]
满足 $20\le\nu_n\le50$kHz 的奇数倍依次是
7、9、11、13，因而
\[
\boxed{4\text{个模式： }
26.6,\ 34.2,\ 41.8,\ 49.4\ \mathrm{kHz}}.
\]
\excheck 使用原版书末近似 $v\simeq3.8$km/s，
模态编号 $n=3,4,5,6$；
原题没有指定铜的物性精确温度；旧稿误把 $E=1.1\times10^{11}\,\mathrm{Pa}$、$\rho=8.9\times10^3\,\mathrm{kg/m^3}$ 代成 $3.8\,\mathrm{km/s}$（该组数据实际约 $3.52\,\mathrm{km/s}$），此次改用与本页频率估计一致的有效模量。

\waveproblem{172}{178}
质量为 $m$ 的两端固定弦作\emph{基频}驻波，
角频率 $\omega$，最大位移振幅为 $a_{\max}$。
求：(a) 最大总动能；
(b) 一个完整周期的平均动能。
\exsolution
取弦长 $L$，
基频位移
$\xi(x,t)=a_{\max}\sin(\pi x/L)\cos\omega t$。
线密度
$\mu=m/L$，
总动能
\[
K(t)=\int_0^L\frac12\mu
(\partial_t\xi)^2\,\mathrm dx
=\frac12\frac mL a_{\max}^2\omega^2
\sin^2\omega t\int_0^L\sin^2(\pi x/L)\,\mathrm dx.
\]
积分等于 $L/2$，故
$K(t)=\frac14m a_{\max}^2\omega^2\sin^2\omega t$。
由此
\[
\boxed{K_{\max}=\frac14m\omega^2a_{\max}^2},
\quad
\boxed{\langle K\rangle=\frac18m\omega^2a_{\max}^2}.
\]
\excheck 平均 $\langle\sin^2\rangle=1/2$，
因而平均动能为最大值的一半。

\waveproblem{173}{178}
均匀长杆横截面积为 $S$、材料密度为 $\rho$，
内部纵向驻波
$\xi=a\sin kx\cos\omega t$。
求相邻两个位移结点之间包含的全部振动机械能。
\exsolution
结点间距为
$\Delta x=\pi/k$。
在 $t=0$ 质点速度处处为零，
机械能全部为弹性势能：
$w_p=\tfrac12\rho v^2(\partial_x\xi)^2
=\tfrac12\rho\omega^2a^2\cos^2kx$，
因为 $v=\omega/k$。
积分：
\[
W=S\int_0^{\pi/k}\frac12
\rho\omega^2a^2\cos^2kx\,\mathrm dx
=\frac12S\rho\omega^2a^2\frac{\pi}{2k}.
\]
因此
\[
\boxed{W=\frac{\pi S\rho\omega^2a^2}{4k}}.
\]
\excheck 对固定的一段相邻结点之间的介质，
无耗散时总机械能保持不变，故可在速度全零的时刻求能。

\waveproblem{174}{178}
频率 $\nu_0=1000$Hz 的声源以
$u=0.17$m/s 沿垂直墙面的直线运动。
两个静止接收器的空间排列为
$R_1$—声源—$R_2$—墙，
声速 $v=340$m/s。
指出哪个接收器会听到拍声以及拍频。
\exsolution
按题设声源向墙一侧移动。
$R_1$ 位于声源后方，接收到直接\emph{远离}声源的频率
$\nu_-= \nu_0v/(v+u)$，
以及在墙上反射、由源\emph{接近}墙面时确定的频率
$\nu_+=\nu_0v/(v-u)$。
因为墙静止，反射频率不变；
所以 $R_1$ 接收两不同频率的波，会产生拍声。
$R_2$ 位于声源前方，其直达与反射波具有相同源趋近频率，不产生这种拍声。
拍频为
\[
\boxed{\nu_b=\nu_+-\nu_-=
\frac{2\nu_0vu}{v^2-u^2}
\simeq\frac{2\nu_0u}{v}
\simeq1.0\,\mathrm{Hz}}.
\]
\excheck 判断接收器的关键是直接声与反射声所对应的源径向运动方向，而不是离墙远近。

\waveproblem{175}{179}
静止观察者同时听到两只频率均为
$\nu_0=680$Hz 的音叉，
一只以速度 $u$ 接近、另一只以同速远离；
听到拍频
$\nu_b=2.0$Hz。
空气声速 $v=340$m/s，求 $u$。
\exsolution
两只运动声源对应接收频率
$\nu_\pm=\nu_0v/(v\mp u)$，
拍频
\[
\nu_b=\nu_+-\nu_-
=\frac{2\nu_0vu}{v^2-u^2}.
\]
化为关于 $u$ 的二次方程，
$\nu_bu^2+2\nu_0vu-\nu_bv^2=0$，
取正根：
\[
\boxed{u=\frac v{\nu_b}
\left(\sqrt{\nu_0^2+\nu_b^2}-\nu_0\right)}.
\]
由于 $\nu_b\ll\nu_0$，
\[
\boxed{u\simeq\frac{v\nu_b}{2\nu_0}
=\frac{340(2)}{2(680)}
\simeq0.50\,\mathrm{m/s}}.
\]
\excheck 用数值计算直接相减两个相近大数容易舍入，
实际可采用等价表达式
$u=v\nu_b/(\sqrt{\nu_0^2+\nu_b^2}+\nu_0)$。

\waveproblem{176}{179}
频率 $\nu_0=2000$Hz 的声源沿接收者与声源连线作简谐往复运动，振幅 $a=50$cm、角频率为 $\Omega$；
静止接收者接收频率的上下限之差
$\Delta\nu=200$Hz，
声速 $v=340$m/s。
求 $\Omega$。
\exsolution
声源沿轴瞬时速度的极大模为
$u_{\max}=a\Omega$。
最高与最低接收频率分别为
$\nu_{\rm hi}=\nu_0v/(v-u_{\max})$、
$\nu_{\rm lo}=\nu_0v/(v+u_{\max})$。
故
\[
\Delta\nu=
\frac{2\nu_0v u_{\max}}{v^2-u_{\max}^2}.
\]
解出
$u_{\max}=v[\sqrt{\nu_0^2+\Delta\nu^2}-\nu_0]/\Delta\nu$，
进而
\[
\boxed{\Omega=\frac v{a\Delta\nu}
\left(\sqrt{\nu_0^2+\Delta\nu^2}-\nu_0\right)}.
\]
取 $\Delta\nu\ll\nu_0$ 近似：
\[
\boxed{\Omega\simeq\frac{v\Delta\nu}
{2a\nu_0}\simeq34\,\mathrm{s^{-1}}}.
\]
\excheck 瞬时频移由声源\emph{速度}决定，
不由位移振幅本身直接决定。

\waveproblem{177}{179}
频率 $\nu_0=1700$Hz 的声源与静止接收者最初位于同一点，
$t=0$ 起声源以
$w=10.0\,\mathrm{m/s^2}$ 的恒加速度远离接收者。
声速 $v=340$m/s。
求从开始运动算起 $t=10.0$s 时接收者\emph{真正听到}的频率。
\exsolution
接收时刻 $t$ 的声波在较早的发声时刻
$t_e$ 发出。
当时声源与接收者相距
$x_e=\tfrac12wt_e^2$，
声波传播历时 $x_e/v$，因此
\[
t=t_e+\frac{wt_e^2}{2v}.
\]
取物理解
$t_e=(v/w)[\sqrt{1+2wt/v}-1]$。
发声时刻的声源速度
$u_e=wt_e=v[\sqrt{1+2wt/v}-1]$。
接收者所听频率为
$\nu=\nu_0v/(v+u_e)$，故
\[
\boxed{\nu(t)=
\frac{\nu_0}{\sqrt{1+2wt/v}}}.
\]
代 $t=10$s，
\[
\boxed{\nu\simeq
\frac{1700}{\sqrt{1+200/340}}
\simeq1.35\,\mathrm{kHz}}.
\]
\excheck 若错误地使用\emph{接收时刻}的声源速度 $wt$，
而不考虑声波传播的迟延，会得到不同答案；
本题必须用发声时刻的速度。

\waveproblem{178}{179}
频率 $\nu_0=1.8$kHz 的点声源沿直线匀速运动，
其运动轨迹到静止观察者的最短距离为
$\ell=250$m，
声源速度 $u=\eta v$，
$\eta=0.80$，$v$ 为声速。
求：(a) 声源\emph{到达最近点的那个瞬间}观察者接收的频率；
(b) 观察者\emph{收到原频率} $\nu_0$ 时，声源与观察者的实际距离。
\exsolution
(a) 设声源最近点时刻 $t=0$，观察者位于 $(0,\ell)$；
此时收到的声波由较早的声源位置
$x_e=-ut_r$ 发出。
传播距离
$r_e=\sqrt{\ell^2+u^2t_r^2}=vt_r$，
解得
$t_r=\ell/[v\sqrt{1-\eta^2}]$。
传播方向上声源的趋近速度投影为
$u^2/v$，
因此接收频率
\[
\boxed{\nu_a=\frac{\nu_0v}{v-u^2/v}
=\frac{\nu_0}{1-\eta^2}
=\frac{1.8}{0.36}\,\mathrm{kHz}
=5.0\,\mathrm{kHz}}.
\]
(b) 要接收到原频率，必须对应\emph{声源恰经过最近点时}发出的声波；
该波抵达观察者需要 $\ell/v$ 时间。
其间声源又沿轨迹移动
$u\ell/v=\eta\ell$，
故接收瞬间的实际距离
\[
\boxed{r_b=\ell\sqrt{1+\eta^2}
=250\sqrt{1.64}\simeq0.32\,\mathrm{km}}.
\]
\excheck 本题(a)需考虑\emph{迟延发声位置}，
不可把当前瞬间径向速度零直接用于多普勒公式；
(b)源当前距离与发射时距离是不同的物理量。

\waveproblem{179}{179}
静止声源发出单色声波，
一面墙以速度
$u=33\,\mathrm{cm/s}$ 朝声源移动。
介质声速
$v=330\,\mathrm{m/s}$。
求反射声的波长相对入射波如何改变，百分比多少。
\exsolution
墙作为运动接收者先听到频率
$\nu_{\rm wall}=\nu_0(v+u)/v$；
反射时墙作为朝声源移动的运动发射者，
其反射波在介质中的频率
$\nu_{\rm ref}=\nu_{\rm wall}v/(v-u)
=\nu_0(v+u)/(v-u)$。
故波长比
\[
\frac{\lambda_{\rm ref}}{\lambda_0}
=\frac{\nu_0}{\nu_{\rm ref}}
=\frac{v-u}{v+u}.
\]
相对减少百分率
\[
\boxed{100\frac{\lambda_0-\lambda_{\rm ref}}
{\lambda_0}=100\frac{2u}{v+u}}
=\boxed{0.20\%}.
\]
\excheck $u=33$cm/s $=0.33$m/s，
所以 $2u/(v+u)\simeq0.0020$，
对应百分数\emph{0.20\%}；
本版书末似印2.0\%，有十倍出入，已记录。

\waveproblem{180}{179}
频率 $\nu_0=1700$Hz 的静止声源与静止接收者位于垂直于一面墙的直线上；
墙以速度
$u=6.0\,\mathrm{cm/s}$ 远离声源。
声速 $v=340$m/s。
求接收者听到的直达声与反射声之间的拍频。
\exsolution
墙作为向外远离的移动接收者，
接收频率
$\nu_{\rm wall}=\nu_0(v-u)/v$。
反射后墙相当于远离静止接收者的移动声源，
反射频率
\[
\nu_{\rm ref}=\frac{\nu_{\rm wall}v}{v+u}
=\nu_0\frac{v-u}{v+u}.
\]
直达频率仍为 $\nu_0$，
所以拍频
\[
\boxed{\nu_b=\nu_0-\nu_{\rm ref}
=\frac{2\nu_0u}{v+u}}.
\]
代 $u=0.060$m/s，
\[
\boxed{\nu_b\simeq0.60\,\mathrm{Hz}}.
\]
\excheck 对运动反射墙要连续应用两次多普勒关系，不能只做一次声源或接收者变换。

\waveproblem{181}{179}
各向同性点声源发声，在距声源
$r_1=10$m、$r_2=20$m 的两点，
声强之比为
$\eta=4.5$。
求声波振幅按
$e^{-\gamma r}$ 衰减时的衰减系数 $\gamma$。
\exsolution
球面波因几何扩散，位移幅度
$A(r)\propto e^{-\gamma r}/r$。
声强与幅度平方成比例，
故
$I(r)\propto e^{-2\gamma r}/r^2$。
因此
\[
\frac{I(r_1)}{I(r_2)}
=\left(\frac{r_2}{r_1}\right)^2
e^{2\gamma(r_2-r_1)}
=\eta.
\]
解出
\[
\boxed{\gamma=\frac1{2(r_2-r_1)}
\ln\left(\eta\frac{r_1^2}{r_2^2}\right)}.
\]
数据：
$\eta r_1^2/r_2^2=4.5/4=1.125$，
所以
\[
\boxed{\gamma=\frac{\ln1.125}{20\,\mathrm m}
\simeq5.9\times10^{-3}\,\mathrm{m^{-1}}
\approx0.006\,\mathrm{m^{-1}}}.
\]
\excheck 若忽略球面扩散，直接用声强比拟合指数衰减，
会高估介质本身的衰减系数。
% === END 25-elastic-waves-166-181.tex ===

% === BEGIN 25-elastic-waves-182-188.tex blob 7cf9d38c59add02798ed3e427f4ac9ccf64828f9 ===
% Irodov Part IV §4.3 182--188, completion of elastic waves and acoustics.
% User source PDF 188--190, printed 179--180; book-end answers 323.
\waveproblem{182}{179--180}
沿 $+x$ 方向传播的平面声波振幅衰减系数 $\gamma=0.0230\,\mathrm{m^{-1}}$；在 $x=0$ 处响度级为 $L=60\,\mathrm{dB}$。求：(a) $x=50\,\mathrm m$ 处的响度级；(b) 声音恰低于听阈的位置。
\exsolution
位移振幅随距离变化 $a(x)=a_0e^{-\gamma x}$；声强与振幅平方成正比，因此 $I(x)=I(0)e^{-2\gamma x}$。
分贝定义 $L(x)=10\log_{10}[I(x)/I_{\rm th}]$ 给
\[
L(x)=L(0)-\frac{20\gamma x}{\ln10}.
\]
(a)
\[
\boxed{L(50)=60-\frac{20(0.0230)(50)}{\ln10}
\simeq50\,\mathrm{dB}}.
\]
(b) 以 $L(x_*)=0$ 为听阈：
\[
\boxed{x_*=\frac{L(0)\ln10}{20\gamma}
\simeq300\,\mathrm m=0.30\,\mathrm{km}}.
\]
\excheck 衰减的是\emph{振幅}，声强的指数为 $-2\gamma x$，不能漏因子2。

\waveproblem{183}{180}
各向同性点声源在 $r_0=20.0\,\mathrm m$ 处响度级 $L_0=30.0\,\mathrm{dB}$。忽略介质吸收，求：(a) $r=10.0\,\mathrm m$ 处响度级；(b) 从多远开始不可听见。
\exsolution
点源球面扩散使 $I(r)\propto r^{-2}$，故
\[
L(r)-L_0=10\log_{10}\frac{I(r)}{I(r_0)}
=20\log_{10}\frac{r_0}{r}.
\]
(a)
$\boxed{L(10)=30+20\log_{10}2
\simeq36\,\mathrm{dB}}$。
(b) 听阈为零分贝，条件
$0=L_0+20\log_{10}(r_0/r_*)$，故
\[
\boxed{r_*=r_0\,10^{L_0/20}
=20\,10^{1.5}\,\mathrm m
\simeq0.63\,\mathrm{km}}.
\]
\excheck 这里分贝基准是相同频率的听阈，且忽略环境噪声与吸收；物理声强并不会在阈值位置突然变成零。

\waveproblem{184}{180}
两个静止观察者 A、B 到正在衰减发声的音叉的距离分别为 $r_A=5.0\,\mathrm m$、$r_B=50\,\mathrm m$。A 听到声音消失的时刻比 B 晚 $\tau=19\,\mathrm s$。声速 $v=340\,\mathrm{m/s}$。求音叉的振幅阻尼系数 $\beta$。
\exsolution
令音叉发声振幅 $a_s(t_e)=a_0e^{-\beta t_e}$。
声从源到半径 $r$ 的接收时刻 $t=t_e+r/v$，接收声压振幅正比于
$e^{-\beta(t-r/v)}/r$。
设听觉阈值声压相同；在各自最后可听时刻 $t_A,t_B$：
\[
\frac{e^{-\beta(t_A-r_A/v)}}{r_A}
=\frac{e^{-\beta(t_B-r_B/v)}}{r_B}.
\]
取对数，使用 $t_A-t_B=\tau$：
$\beta[\tau+(r_B-r_A)/v]=\ln(r_B/r_A)$。
因此
\[
\boxed{\beta=\frac{\ln(r_B/r_A)}
{\tau+(r_B-r_A)/v}
=\frac{\ln10}{19+45/340}
\simeq0.12\,\mathrm{s^{-1}}}.
\]
\excheck 要考虑额外 $45$m 的声传播迟延，而非仅用 $\ln10/19$；影响虽小但原题要求保留。

\waveproblem{185}{180}
密度为 $\rho$ 的均匀介质中传播平面纵向谐波，传播速度 $v$，扰动密度很小。证明：(a) 声压增量 $\Delta p=-\rho v^2\,\partial_x\xi$；(b) 平均声强与波振幅和能量密度的关系。
\exsolution
对截面积 $S$、厚度 $\mathrm dx$ 的薄介质层用牛顿第二定律，
$\rho S\mathrm dx\,\partial_t^2\xi=-S(\partial_xp)\mathrm dx$，
即
$\rho\partial_t^2\xi=-\partial_x\Delta p$。
由一维波动方程
$\partial_t^2\xi=v^2\partial_x^2\xi$，
得
$\partial_x[\Delta p+\rho v^2\partial_x\xi]=0$。
无扰动时附加压力与应变均为零，积分常数为零，
\[
\boxed{\Delta p=-\rho v^2\frac{\partial\xi}{\partial x}}.
\]
对于行波 $\xi=a\cos(\omega t-kx)$，
$p'= -\rho v^2ak\sin(\omega t-kx)$，
粒子速度 $u=-a\omega\sin(\omega t-kx)$；
传播方向的瞬时能流密度 $S=p'u$，
故
\[
\boxed{\overline S
=\tfrac12\rho v\,\omega^2a^2
=\frac{(\Delta p)_m^2}{2\rho v}
=\overline w\,v}.
\]
\excheck 压强扰动与粒子速度同相，使平均声功率沿行波传播方向为正；若波朝反方向行进，能流矢量亦反向。

\waveproblem{186}{180}
半径 $R=50\,\mathrm{cm}$ 的球置于平面声波路径上，声波波长 $\lambda=20$cm、频率 $\nu=1700$Hz、声压振幅 $(\Delta p)_m=3.5$Pa。求一个周期平均抵达球表面的能流功率。
\exsolution
空气声速 $v=\nu\lambda=340$m/s。
由第185题平均声强
$\overline I=(\Delta p)_m^2/(2\rho v)$，
球截获迎波有效截面积为 $\pi R^2$（不是完整球表面积 $4\pi R^2$），所以
\[
\boxed{\Phi=\pi R^2\,\frac{(\Delta p)_m^2}{2\rho\nu\lambda}}.
\]
以空气密度 $\rho\simeq1.29\,\mathrm{kg/m^3}$ 估算：
\[
\boxed{\Phi\simeq
\frac{\pi(0.50)^2(3.5)^2}{2(1.29)(340)}
\simeq0.011\,\mathrm W=11\,\mathrm{mW}}.
\]
\excheck 平面行波对球面的法向能流积分等于投影圆截面乘声强；即使球面面积更大也不能直接乘 $4\pi R^2$。

\waveproblem{187}{180}
点声源频率 $\nu=600$Hz、声功率 $P=0.80$W，在距声源 $r=1.5$m 的 A 点，取声速 $v=340$m/s，忽略衰减。求：(a) 声压振幅及其与大气压之比；(b) 介质质点位移幅度及它与波长之比。
\exsolution
各向同性点声源在 A 点的平均声强
$\overline I=P/(4\pi r^2)$。
由第185题
$\overline I=(\Delta p)_m^2/(2\rho v)$，故
\[
\boxed{(\Delta p)_m=
\sqrt{\frac{\rho vP}{2\pi r^2}}
\simeq5\,\mathrm{Pa}}.
\]
取大气压 $p_{\rm atm}\simeq10^5$Pa，
$\boxed{(\Delta p)_m/p_{\rm atm}\simeq5\times10^{-5}}$。
声压—位移关系 $(\Delta p)_m=\rho v\omega a$，
$\omega=2\pi\nu$，所以
\[
\boxed{a=\frac{(\Delta p)_m}{2\pi\nu\rho v}
\simeq3\,\mathrm{\mu m}},
\qquad
\boxed{\frac a\lambda=\frac{a\nu}{v}
\simeq5\times10^{-6}}.
\]
\excheck 声传播速度 $v$ 与介质质点速度幅度 $\omega a$ 是不同量；本题声压扰动相对大气压极小，线性声学适用。

\waveproblem{188}{180}
距各向同性点声源 $r=100$m 处，频率为200Hz的声音响度级 $L=50$dB；该频率的听阈声强 $I_0=0.10\,\mathrm{nW/m^2}$。振幅衰减系数 $\gamma=5.0\times10^{-4}\,\mathrm{m^{-1}}$。求点声源功率。
\exsolution
声强定义
$I(r)=I_0\,10^{L/10}
=10^{-10}\times10^5=10^{-5}\,\mathrm{W/m^2}$。
球面波幅度为 $a\propto e^{-\gamma r}/r$，
故声强为 $I=P e^{-2\gamma r}/(4\pi r^2)$。
解得
\[
\boxed{P=4\pi r^2 I_0\,10^{L/10}e^{2\gamma r}}.
\]
代入 $2\gamma r=0.10$：
\[
\boxed{P=4\pi(100)^2(10^{-5})e^{0.10}
\simeq1.4\,\mathrm W}.
\]
\excheck $I_0$ 为 $0.10$纳瓦每平方米，必须换算为 $10^{-10}\,\mathrm{W/m^2}$；分贝 $50$dB 对应强度比 $10^5$ 而非50。
% === END 25-elastic-waves-182-188.tex ===
